\section{Explain the measured work at a complete function boundary}
The public eight-pixel trace occupies thirty-eight bytes. Its explicit red command is
\texttt{fe ff 00 00}: four bytes to replace RGB while retaining alpha. The preceding
pixel is opaque black. Red can also be obtained with changes \((-1,0,0)\), wrapping
zero minus one to 255. Biasing those three DIFF fields by two gives \texttt{01 10 10};
with the class bits, the complete tag is \texttt{5a}. That single byte saves three
bytes without changing the reconstructed color or cache state.

The reference prototype chooses that shorter command and writes a thirty-five-byte
file for the trace. Its decoded pixels and descriptive header fields match the
thirty-eight-byte fixture. The comparison is between two representations of the same
image, so the encoder check decodes the result independently rather than comparing
file bytes with the fixture. For the sprite, the reference produces sixty-one bytes.

The following measurements count completed simulator instructions from the named
function entry through its return, including private cache initialization, called
helpers, validation and delayed returns. Preparing arguments and selecting the stopped
entry externally are not executed program instructions. Validate and decode receive
the same public fixture bytes; encode receives the corresponding RGBA words and
metadata. The last column is the encoder's returned file length.

\tagpdfsetup{table/header-rows={1}}
\begin{center}
  \small
  \begin{tabular}{llrrrr}
    Fixture           & Prototype   & Validate & Decode & Encode & File bytes \\
    \hline
    Eight-pixel trace & Reference   & 846      & 1748   & 772    & 35         \\
    Eight-pixel trace & Alternative & 854      & 1764   & 788    & 35         \\
    Eight-row sprite  & Reference   & 2690     & 5548   & 2813   & 61         \\
    Eight-row sprite  & Alternative & 2766     & 5700   & 2941   & 61         \\
  \end{tabular}
\end{center}

The decoder performs validation and a second reconstruction pass, plus its own capacity,
frame and return work. Its count therefore exceeds twice the standalone validation
count in these examples. This is the measured cost of retaining the unchanged-output
failure promise. For these prototypes, the difference is two additional instructions
per pixel and forty fixed instructions: the output store and pointer advance add
\(2n\), and the wrapper adds forty. Thus the trace has \(1748-2\cdot846=56\) extra
instructions, while the sprite has \(5548-2\cdot2690=168\). This relationship is an
observation of these implementations, not a universal decoder formula. Removing the
validation pass would alter the failure promise, not merely remove redundant instructions
from the same interface.

The alternative uses bit-field insertion when assembling colors and shifts and additions
for the cache hash. It also permits a different literal choice during encoding. These
are valid implementation choices, but the measured alternative takes more instructions
for these two inputs. Replacing multiplication with several shifts and additions does
not automatically reduce simulated work. A host-time comparison would need a separate
measurement; this table does not measure wall-clock execution or IDE launch time.

Treat these counts as observations of the prototypes, not required answers for a
student implementation. A different correct cache representation, register allocation
or encoder choice can produce other counts. Compare complete boundaries on matching
inputs before attributing a change to a particular instruction sequence.
